% Part: computability
% Chapter: computability-theory
% Section: ce-sets

\documentclass[../../../include/open-logic-section]{subfiles}

\begin{document}

\olfileid{cmp}{thy}{ces}
\olsection{可计算枚举集}

\begin{defn}
一个集合是\emph{可计算枚举的}，如果它是空集，或是某个可计算函数的值域。
\end{defn}

\begin{history}
可计算枚举集也可称为\emph{递归可枚举}集。这是最初的术语，如今两者及其缩写“c.e.”和“r.e.”都很常用。
\end{history}

\begin{explain}
请思考这一定义的含义，以及为何该术语是恰当的。其直观是，如果 $S$ 是可计算函数~$f$ 的值域，那么
\[
S = \{ f(0), f(1), f(2), \dots \},
\]
因此 $f$ 可以看作“枚举”了~$S$ 的元素。注意，根据定义，$f$~不必是递增函数，即枚举的顺序不必是递增的。事实上，$f$ 甚至不必是单射的，即在~$S$ 的枚举 $f(0)$、$f(1)$、$f(2)$、\dots{} 中允许出现重复项。例如，常函数 $f(x) = 0$ 就枚举了集合 $\{ 0 \}$。
\end{explain}

任何可计算集都是可计算枚举的。为说明这一点，设 $S$~是可计算的。如果 $S$ 是空集，那么根据定义它是可计算枚举的。否则，令 $a$ 为 $S$ 中的任意元素。按下式定义 $f$：
\[
f(x) =
\begin{cases}
x & \text{如果 $\Char{S}(x) = 1$} \\
a & \text{否则。}
\end{cases}
\]
那么 $f$ 是可计算函数，且 $S$ 是~$f$ 的值域。

\end{document}